\section{Group actions}\label{Groups}
In this section we recall basic properties of the $R$-tensor product $\otimes_R$ and the ring dual $\Hom_R(-,R)$ of $R$-modules. Both constructions extend to chain complexes of $R$-modules by applying them degreewise.

\subsection{The $R$-tensor}

Let $R$ be a ring, let $M$ be a right $R$-module, and let $N$ be a left $R$-module. The \defn{$R$-tensor product} $M \otimes_R N$ is the quotient
\[
M \otimes_R N \coloneqq (M \otimes N) \big/ (mr \otimes n \sim m \otimes rn).
\]
As defined, $M \otimes_R N$ is an abelian group, and does not carry a natural $R$-module structure in general.

It is often convenient to view $\otimes_R$ as a monoidal product in a category of bimodules whose bjects are triples $(R,S,N)$ where $R$ and $S$ are rings and $N$ is an $(R,S)$-bimodule. Given objects $(T,R,M)$ and $(R,S,N)$, their monoidal product is $(T,S,M \otimes_R N)$, where $M \otimes_R N$ is a $(T,S)$-bimodule with left $T$-action and right $S$-action given by
\[
t \cdot (m \otimes n) = tm \otimes n,
\qquad
(m \otimes n) \cdot s = m \otimes ns.
\]

If $M$ is instead a left $R$-module, additional structure is required in order to define $M \otimes_R N$. We record two important cases.

\begin{example}
Suppose $R$ is commutative. Any left $R$-module $M$ acquires a right $R$-module structure by setting $m \cdot r \coloneqq rm$. In this case, the tensor product is given by
\[
M \otimes_R N = (M \otimes N) \big/ (rm \otimes n \sim m \otimes rn).
\]
This defines the monoidal product on the category of $R$-modules, which is a symmetric monoidal category and can be seen as the restriction of the category of bimodules to the full subcategory whose objects are triplets $(R,R, M)$ in which the left and right $R$-actions on $M$ coincide. The tensor product $M \otimes_R N$ is again an $R$-module with action
\[
r \cdot (m \otimes n) = rm \otimes n.
\]

The category of $R$-modules also admits an internal hom, given by $\Hom_R$, with $R$-module structure defined by
\[
(r \cdot f)(m) = f(rm) = r f(m).
\]
Moreover, for $R$-modules $M$, $N$, and $P$, there is a natural isomorphism
\[
\Hom_R(M \otimes_R N, P) \cong \Hom_R(M, \Hom_R(N,P)),
\]
known as the \defn{hom--tensor adjunction}.
\end{example}

\begin{example}
Let $R = \k[G]$ for a field $\k$ and a group $G$. Any left $\k[G]$-module $M$ can be endowed with a right $\k[G]$-module structure via the group anti-automorphism $g \mapsto g^{-1}$, by defining
\[
m \cdot g \coloneqq g^{-1} \cdot m.
\]
This does not in general define a bimodule structure unless $G$ is abelian.

We denote the resulting tensor product by $\otimes_G$, and define
\[
M \otimes_G N \coloneqq (M \otimes N) \big/ (g^{-1}m \otimes n \sim m \otimes gn).
\]
Replacing $m$ by $gm$ in the defining relation shows that $(m \otimes n) \sim (gm \otimes gn)$, and hence
\[
M \otimes_G N \cong (M \otimes N)_G \coloneqq \k \otimes_{\k[G]} (M \otimes N),
\]
the \defn{coinvariants} of $M \otimes N$ with respect to the diagonal $G$-action
\[
g \cdot (m \otimes n) = gm \otimes gn.
\]

There is also a diagonal $G$-action on $\Hom(M,N)$ given by
\[
(g \cdot f)(m) = g f(g^{-1} m),
\]
which identifies equivariant maps with \defn{invariants}:
\[
\Hom_G(M,N) = (\Hom(M,N))^G \coloneqq \Hom_{\k[G]}(\k, \Hom(M,N)).
\]
Note that, by definition, $\Hom_G(M,N) = \Hom_{\k[G]}(M,N)$.

If $A$ is a $\k$-vector space with the trivial $G$-action, then there are natural identifications
\[
\Hom(M_G, A) \cong \Hom_G(M,A),
\qquad
\Hom_G(A,M) \cong \Hom(A, M^G),
\]
relating invariants and coinvariants.
\end{example}

\subsection{The ring dual}

Let $R$ be a ring and let $M$ be a left $R$-module. The \defn{ring dual} of $M$ is the right $R$-module
\[
M^\vee \coloneqq \Hom_R(M,R),
\]
with right $R$-action given by
\[
(f \cdot r)(m) = f(m) r.
\]

If $M$ is finitely generated and projective as a left $R$-module, then $M^\vee$ is finitely generated and projective as a right $R$-module. Moreover, for any left $R$-module $N$, there is an isomorphism
\[
\varphi \colon M^\vee \otimes_R N \xrightarrow{\sim} \Hom_R(M,N),
\]
defined by
\[
\varphi(f \otimes n)(m) = f(m) n.
\]

\begin{example}\label{BimoduleW}
Let $R = \k[G]$ for a field $\k$ and a group $G$, and consider the left $R$-module $W_0 \coloneqq R\{e_0\}$ with the regular action $g \cdot e_0 = ge_0$. Equivalently, we may identify $W_0$ with $\k\{ge_0 \mid g \in G\}$, where $G$ acts by permuting the basis.

The dual module is
\[
W_0^\vee = \Hom_R(W_0,R) = \{e_0^\vee\}R = \k\{(ge_0)^\vee \mid g \in G\},
\]
where $(ge_0)^\vee$ denotes the $R$-linear map sending $ge_0$ to $1_R$. The right $R$-action is given by
\[
e_0^\vee \cdot g = (g^{-1} e_0)^\vee.
\]

Recall that $W_0$ admits a right $R$-module structure defined by $e_0 \cdot g = g^{-1} e_0$. If $G$ is abelian, then $W_0$ becomes a bimodule, and this induces a left $R$-action on $W_0^\vee$ given by
\[
(g \cdot e_0^\vee)(he_0) \coloneqq e_0^\vee(he_0 \cdot g)= (ge_0)^\vee(he_0).
\]
This is well-defined precisely because the bimodule condition holds.

Finally, note that the right action on $W_0^\vee$ can equivalently be described as
\[
e_0^\vee \cdot g = g^{-1} \cdot e_0^\vee,
\]
which agrees with the right module structure used in the construction of the tensor product $W_0^\vee \otimes_G -$.
\end{example}
